% Priority inversion and priority inheritance.  T1 (high) arrives at 3,
% T2 (medium) at 2, T3 (low) at 0; T1 and T3 share mutex m.
% Usage: % Priority inversion and priority inheritance.  T1 (high) arrives at 3,
% T2 (medium) at 2, T3 (low) at 0; T1 and T3 share mutex m.
% Usage: % Priority inversion and priority inheritance.  T1 (high) arrives at 3,
% T2 (medium) at 2, T3 (low) at 0; T1 and T3 share mutex m.
% Usage: % Priority inversion and priority inheritance.  T1 (high) arrives at 3,
% T2 (medium) at 2, T3 (low) at 0; T1 and T3 share mutex m.
% Usage: \input{figures/l07-inversion}   (width ~112 mm)
\begin{tikzpicture}[x=9mm, y=1mm, font=\footnotesize\sffamily,
  run/.style={draw, fill=ln-box},
  cs/.style={draw, fill=ln-accent!40},
  boost/.style={draw=ln-caution, line width=1.2pt, fill=ln-accent!40},
  blk/.style={densely dashed, thick, ln-caution},
  lab/.style={anchor=east},
  ttl/.style={anchor=south west, font=\footnotesize\sffamily\bfseries, inner sep=1pt}]
  % ---------- (a) without priority inheritance
  \node[ttl] at (0,54) {\iflangja{(a) 優先度継承なし}{(a) without priority inheritance}};
  \node[lab] at (-0.1,50.5) {\iflangja{T1（高）}{T1 (high)}};
  \node[lab] at (-0.1,42.5) {\iflangja{T2（中）}{T2 (medium)}};
  \node[lab] at (-0.1,34.5) {\iflangja{T3（低）}{T3 (low)}};
  \draw[run] (3,48) rectangle (4,53);
  \draw[blk] (4,50.5) -- (7,50.5);
  \draw[cs]  (7,48) rectangle (8,53);
  \draw[run] (8,48) rectangle (9,53);
  \draw[run] (2,40) rectangle (3,45);
  \draw[run] (4,40) rectangle (6,45);
  \draw[run] (0,32) rectangle (1,37);
  \draw[cs]  (1,32) rectangle (2,37);
  \draw[cs]  (6,32) rectangle (7,37);
  \draw[run] (9,32) rectangle (10,37);
  \draw (0,29) -- (10.3,29);
  \foreach \x in {0,...,10} {
    \draw (\x,29) -- (\x,28);
    \node[below, font=\scriptsize\sffamily, inner sep=1pt] at (\x,28) {\x};
  }
  % ---------- (b) with priority inheritance
  \node[ttl] at (0,19) {\iflangja{(b) 優先度継承あり}{(b) with priority inheritance}};
  \node[lab] at (-0.1,15.5) {\iflangja{T1（高）}{T1 (high)}};
  \node[lab] at (-0.1,7.5)  {\iflangja{T2（中）}{T2 (medium)}};
  \node[lab] at (-0.1,-0.5) {\iflangja{T3（低）}{T3 (low)}};
  \draw[run] (3,13) rectangle (4,18);
  \draw[blk] (4,15.5) -- (5,15.5);
  \draw[cs]  (5,13) rectangle (6,18);
  \draw[run] (6,13) rectangle (7,18);
  \draw[run] (2,5) rectangle (3,10);
  \draw[run] (7,5) rectangle (9,10);
  \draw[run] (0,-3) rectangle (1,2);
  \draw[cs]  (1,-3) rectangle (2,2);
  \draw[boost] (4,-3) rectangle (5,2);
  \node[anchor=west, font=\scriptsize\sffamily, text=ln-caution] at (5.1,-0.5)
    {\iflangja{優先度を「高」に引き上げ}{priority raised to high}};
  \draw[run] (9,-3) rectangle (10,2);
  \draw (0,-6) -- (10.3,-6);
  \foreach \x in {0,...,10} {
    \draw (\x,-6) -- (\x,-7);
    \node[below, font=\scriptsize\sffamily, inner sep=1pt] at (\x,-7) {\x};
  }
  \node[anchor=west, font=\scriptsize\sffamily] at (10.3,-6) {$t$};
  % ---------- legend
  \begin{scope}[yshift=-17mm, font=\scriptsize\sffamily]
    \draw[run] (0,-1.5) rectangle (0.5,1.5);
    \node[anchor=west] at (0.55,0) {\iflangja{実行}{runs}};
    \draw[cs] (2,-1.5) rectangle (2.5,1.5);
    \node[anchor=west] at (2.55,0) {\iflangja{m を保持して実行}{runs holding m}};
    % The Japanese labels are wider: shift the last two entries so that the
    % "inherited priority" swatch does not cover the end of "blocked on m".
    \draw[blk] (\iflangja{5.0}{5.3},0) -- (\iflangja{5.6}{6},0);
    \node[anchor=west] at (\iflangja{5.65}{6.05},0) {\iflangja{m 待ちでブロック}{blocked on m}};
    \draw[boost] (\iflangja{8.45}{8.2},-1.5) rectangle (\iflangja{8.95}{8.7},1.5);
    \node[anchor=west] at (\iflangja{9.0}{8.75},0) {\iflangja{継承した優先度}{inherited priority}};
  \end{scope}
\end{tikzpicture}
   (width ~112 mm)
\begin{tikzpicture}[x=9mm, y=1mm, font=\footnotesize\sffamily,
  run/.style={draw, fill=ln-box},
  cs/.style={draw, fill=ln-accent!40},
  boost/.style={draw=ln-caution, line width=1.2pt, fill=ln-accent!40},
  blk/.style={densely dashed, thick, ln-caution},
  lab/.style={anchor=east},
  ttl/.style={anchor=south west, font=\footnotesize\sffamily\bfseries, inner sep=1pt}]
  % ---------- (a) without priority inheritance
  \node[ttl] at (0,54) {\iflangja{(a) 優先度継承なし}{(a) without priority inheritance}};
  \node[lab] at (-0.1,50.5) {\iflangja{T1（高）}{T1 (high)}};
  \node[lab] at (-0.1,42.5) {\iflangja{T2（中）}{T2 (medium)}};
  \node[lab] at (-0.1,34.5) {\iflangja{T3（低）}{T3 (low)}};
  \draw[run] (3,48) rectangle (4,53);
  \draw[blk] (4,50.5) -- (7,50.5);
  \draw[cs]  (7,48) rectangle (8,53);
  \draw[run] (8,48) rectangle (9,53);
  \draw[run] (2,40) rectangle (3,45);
  \draw[run] (4,40) rectangle (6,45);
  \draw[run] (0,32) rectangle (1,37);
  \draw[cs]  (1,32) rectangle (2,37);
  \draw[cs]  (6,32) rectangle (7,37);
  \draw[run] (9,32) rectangle (10,37);
  \draw (0,29) -- (10.3,29);
  \foreach \x in {0,...,10} {
    \draw (\x,29) -- (\x,28);
    \node[below, font=\scriptsize\sffamily, inner sep=1pt] at (\x,28) {\x};
  }
  % ---------- (b) with priority inheritance
  \node[ttl] at (0,19) {\iflangja{(b) 優先度継承あり}{(b) with priority inheritance}};
  \node[lab] at (-0.1,15.5) {\iflangja{T1（高）}{T1 (high)}};
  \node[lab] at (-0.1,7.5)  {\iflangja{T2（中）}{T2 (medium)}};
  \node[lab] at (-0.1,-0.5) {\iflangja{T3（低）}{T3 (low)}};
  \draw[run] (3,13) rectangle (4,18);
  \draw[blk] (4,15.5) -- (5,15.5);
  \draw[cs]  (5,13) rectangle (6,18);
  \draw[run] (6,13) rectangle (7,18);
  \draw[run] (2,5) rectangle (3,10);
  \draw[run] (7,5) rectangle (9,10);
  \draw[run] (0,-3) rectangle (1,2);
  \draw[cs]  (1,-3) rectangle (2,2);
  \draw[boost] (4,-3) rectangle (5,2);
  \node[anchor=west, font=\scriptsize\sffamily, text=ln-caution] at (5.1,-0.5)
    {\iflangja{優先度を「高」に引き上げ}{priority raised to high}};
  \draw[run] (9,-3) rectangle (10,2);
  \draw (0,-6) -- (10.3,-6);
  \foreach \x in {0,...,10} {
    \draw (\x,-6) -- (\x,-7);
    \node[below, font=\scriptsize\sffamily, inner sep=1pt] at (\x,-7) {\x};
  }
  \node[anchor=west, font=\scriptsize\sffamily] at (10.3,-6) {$t$};
  % ---------- legend
  \begin{scope}[yshift=-17mm, font=\scriptsize\sffamily]
    \draw[run] (0,-1.5) rectangle (0.5,1.5);
    \node[anchor=west] at (0.55,0) {\iflangja{実行}{runs}};
    \draw[cs] (2,-1.5) rectangle (2.5,1.5);
    \node[anchor=west] at (2.55,0) {\iflangja{m を保持して実行}{runs holding m}};
    % The Japanese labels are wider: shift the last two entries so that the
    % "inherited priority" swatch does not cover the end of "blocked on m".
    \draw[blk] (\iflangja{5.0}{5.3},0) -- (\iflangja{5.6}{6},0);
    \node[anchor=west] at (\iflangja{5.65}{6.05},0) {\iflangja{m 待ちでブロック}{blocked on m}};
    \draw[boost] (\iflangja{8.45}{8.2},-1.5) rectangle (\iflangja{8.95}{8.7},1.5);
    \node[anchor=west] at (\iflangja{9.0}{8.75},0) {\iflangja{継承した優先度}{inherited priority}};
  \end{scope}
\end{tikzpicture}
   (width ~112 mm)
\begin{tikzpicture}[x=9mm, y=1mm, font=\footnotesize\sffamily,
  run/.style={draw, fill=ln-box},
  cs/.style={draw, fill=ln-accent!40},
  boost/.style={draw=ln-caution, line width=1.2pt, fill=ln-accent!40},
  blk/.style={densely dashed, thick, ln-caution},
  lab/.style={anchor=east},
  ttl/.style={anchor=south west, font=\footnotesize\sffamily\bfseries, inner sep=1pt}]
  % ---------- (a) without priority inheritance
  \node[ttl] at (0,54) {\iflangja{(a) 優先度継承なし}{(a) without priority inheritance}};
  \node[lab] at (-0.1,50.5) {\iflangja{T1（高）}{T1 (high)}};
  \node[lab] at (-0.1,42.5) {\iflangja{T2（中）}{T2 (medium)}};
  \node[lab] at (-0.1,34.5) {\iflangja{T3（低）}{T3 (low)}};
  \draw[run] (3,48) rectangle (4,53);
  \draw[blk] (4,50.5) -- (7,50.5);
  \draw[cs]  (7,48) rectangle (8,53);
  \draw[run] (8,48) rectangle (9,53);
  \draw[run] (2,40) rectangle (3,45);
  \draw[run] (4,40) rectangle (6,45);
  \draw[run] (0,32) rectangle (1,37);
  \draw[cs]  (1,32) rectangle (2,37);
  \draw[cs]  (6,32) rectangle (7,37);
  \draw[run] (9,32) rectangle (10,37);
  \draw (0,29) -- (10.3,29);
  \foreach \x in {0,...,10} {
    \draw (\x,29) -- (\x,28);
    \node[below, font=\scriptsize\sffamily, inner sep=1pt] at (\x,28) {\x};
  }
  % ---------- (b) with priority inheritance
  \node[ttl] at (0,19) {\iflangja{(b) 優先度継承あり}{(b) with priority inheritance}};
  \node[lab] at (-0.1,15.5) {\iflangja{T1（高）}{T1 (high)}};
  \node[lab] at (-0.1,7.5)  {\iflangja{T2（中）}{T2 (medium)}};
  \node[lab] at (-0.1,-0.5) {\iflangja{T3（低）}{T3 (low)}};
  \draw[run] (3,13) rectangle (4,18);
  \draw[blk] (4,15.5) -- (5,15.5);
  \draw[cs]  (5,13) rectangle (6,18);
  \draw[run] (6,13) rectangle (7,18);
  \draw[run] (2,5) rectangle (3,10);
  \draw[run] (7,5) rectangle (9,10);
  \draw[run] (0,-3) rectangle (1,2);
  \draw[cs]  (1,-3) rectangle (2,2);
  \draw[boost] (4,-3) rectangle (5,2);
  \node[anchor=west, font=\scriptsize\sffamily, text=ln-caution] at (5.1,-0.5)
    {\iflangja{優先度を「高」に引き上げ}{priority raised to high}};
  \draw[run] (9,-3) rectangle (10,2);
  \draw (0,-6) -- (10.3,-6);
  \foreach \x in {0,...,10} {
    \draw (\x,-6) -- (\x,-7);
    \node[below, font=\scriptsize\sffamily, inner sep=1pt] at (\x,-7) {\x};
  }
  \node[anchor=west, font=\scriptsize\sffamily] at (10.3,-6) {$t$};
  % ---------- legend
  \begin{scope}[yshift=-17mm, font=\scriptsize\sffamily]
    \draw[run] (0,-1.5) rectangle (0.5,1.5);
    \node[anchor=west] at (0.55,0) {\iflangja{実行}{runs}};
    \draw[cs] (2,-1.5) rectangle (2.5,1.5);
    \node[anchor=west] at (2.55,0) {\iflangja{m を保持して実行}{runs holding m}};
    % The Japanese labels are wider: shift the last two entries so that the
    % "inherited priority" swatch does not cover the end of "blocked on m".
    \draw[blk] (\iflangja{5.0}{5.3},0) -- (\iflangja{5.6}{6},0);
    \node[anchor=west] at (\iflangja{5.65}{6.05},0) {\iflangja{m 待ちでブロック}{blocked on m}};
    \draw[boost] (\iflangja{8.45}{8.2},-1.5) rectangle (\iflangja{8.95}{8.7},1.5);
    \node[anchor=west] at (\iflangja{9.0}{8.75},0) {\iflangja{継承した優先度}{inherited priority}};
  \end{scope}
\end{tikzpicture}
   (width ~112 mm)
\begin{tikzpicture}[x=9mm, y=1mm, font=\footnotesize\sffamily,
  run/.style={draw, fill=ln-box},
  cs/.style={draw, fill=ln-accent!40},
  boost/.style={draw=ln-caution, line width=1.2pt, fill=ln-accent!40},
  blk/.style={densely dashed, thick, ln-caution},
  lab/.style={anchor=east},
  ttl/.style={anchor=south west, font=\footnotesize\sffamily\bfseries, inner sep=1pt}]
  % ---------- (a) without priority inheritance
  \node[ttl] at (0,54) {\iflangja{(a) 優先度継承なし}{(a) without priority inheritance}};
  \node[lab] at (-0.1,50.5) {\iflangja{T1（高）}{T1 (high)}};
  \node[lab] at (-0.1,42.5) {\iflangja{T2（中）}{T2 (medium)}};
  \node[lab] at (-0.1,34.5) {\iflangja{T3（低）}{T3 (low)}};
  \draw[run] (3,48) rectangle (4,53);
  \draw[blk] (4,50.5) -- (7,50.5);
  \draw[cs]  (7,48) rectangle (8,53);
  \draw[run] (8,48) rectangle (9,53);
  \draw[run] (2,40) rectangle (3,45);
  \draw[run] (4,40) rectangle (6,45);
  \draw[run] (0,32) rectangle (1,37);
  \draw[cs]  (1,32) rectangle (2,37);
  \draw[cs]  (6,32) rectangle (7,37);
  \draw[run] (9,32) rectangle (10,37);
  \draw (0,29) -- (10.3,29);
  \foreach \x in {0,...,10} {
    \draw (\x,29) -- (\x,28);
    \node[below, font=\scriptsize\sffamily, inner sep=1pt] at (\x,28) {\x};
  }
  % ---------- (b) with priority inheritance
  \node[ttl] at (0,19) {\iflangja{(b) 優先度継承あり}{(b) with priority inheritance}};
  \node[lab] at (-0.1,15.5) {\iflangja{T1（高）}{T1 (high)}};
  \node[lab] at (-0.1,7.5)  {\iflangja{T2（中）}{T2 (medium)}};
  \node[lab] at (-0.1,-0.5) {\iflangja{T3（低）}{T3 (low)}};
  \draw[run] (3,13) rectangle (4,18);
  \draw[blk] (4,15.5) -- (5,15.5);
  \draw[cs]  (5,13) rectangle (6,18);
  \draw[run] (6,13) rectangle (7,18);
  \draw[run] (2,5) rectangle (3,10);
  \draw[run] (7,5) rectangle (9,10);
  \draw[run] (0,-3) rectangle (1,2);
  \draw[cs]  (1,-3) rectangle (2,2);
  \draw[boost] (4,-3) rectangle (5,2);
  \node[anchor=west, font=\scriptsize\sffamily, text=ln-caution] at (5.1,-0.5)
    {\iflangja{優先度を「高」に引き上げ}{priority raised to high}};
  \draw[run] (9,-3) rectangle (10,2);
  \draw (0,-6) -- (10.3,-6);
  \foreach \x in {0,...,10} {
    \draw (\x,-6) -- (\x,-7);
    \node[below, font=\scriptsize\sffamily, inner sep=1pt] at (\x,-7) {\x};
  }
  \node[anchor=west, font=\scriptsize\sffamily] at (10.3,-6) {$t$};
  % ---------- legend
  \begin{scope}[yshift=-17mm, font=\scriptsize\sffamily]
    \draw[run] (0,-1.5) rectangle (0.5,1.5);
    \node[anchor=west] at (0.55,0) {\iflangja{実行}{runs}};
    \draw[cs] (2,-1.5) rectangle (2.5,1.5);
    \node[anchor=west] at (2.55,0) {\iflangja{m を保持して実行}{runs holding m}};
    % The Japanese labels are wider: shift the last two entries so that the
    % "inherited priority" swatch does not cover the end of "blocked on m".
    \draw[blk] (\iflangja{5.0}{5.3},0) -- (\iflangja{5.6}{6},0);
    \node[anchor=west] at (\iflangja{5.65}{6.05},0) {\iflangja{m 待ちでブロック}{blocked on m}};
    \draw[boost] (\iflangja{8.45}{8.2},-1.5) rectangle (\iflangja{8.95}{8.7},1.5);
    \node[anchor=west] at (\iflangja{9.0}{8.75},0) {\iflangja{継承した優先度}{inherited priority}};
  \end{scope}
\end{tikzpicture}
